\subsection{Abelianization of a semidirect product}
\begin{lemma}[Cited group formula, with a re-derivation]
\label{wb:lem:semidirect-abelianization}
Let $\alpha:Q\to\Aut(K)$ be an action of a group $Q$ on a group $K$,
and form $G=K\rtimes_\alpha Q$. Put
\[
 C=K^{\mathrm{ab}}/
   \langle\alpha(q)[k]-[k]:q\in Q,\ k\in K\rangle.
\]
The inclusion of $K$ and the chosen section $Q\to G$ induce an
isomorphism of abelian groups
\[
 G^{\mathrm{ab}}\simeq Q^{\mathrm{ab}}\oplus C,
 \qquad C=(K^{\mathrm{ab}})_Q.
 \tag{SD.1}\label{wb:eq:semidirect-abelianization}
\]
Thus the assertion applies to any split group extension once a
splitting has been selected.
\end{lemma}
This formula is \citet[Section~10.5, Lemma~10.7, p.~90]{AC2026}.
We retain the quotient argument of the working edition as the following
independently worded re-derivation.
\Uses{The definition of a semidirect product, a commutator subgroup,
and the displayed coinvariant quotient}
\begin{proof}
Write $s:Q\to G$ for the section. Define $N$ to be the normal subgroup
of $G$ generated by
\[
 [k,k']\quad\hbox{and}\quad s(q)ks(q)^{-1}k^{-1}
 \qquad(k,k'\in K, q\in Q).
\]
Every generator lies in $K\cap[G,G]$, and this intersection is normal
in $G$. Consequently $N\subset K\cap[G,G]$.
The relations defining $K/N$ impose commutativity of $K$ and identify
the class of $k$ with that of $\alpha(q)k$. They therefore give the
canonical isomorphism $K/N\simeq C$. More explicitly, the map $K\to C$
has a $G$-invariant kernel and kills all the displayed generators;
in the reverse direction $K\to K/N$ factors through $K^{\mathrm{ab}}$
and through its $Q$-coinvariants. These two factorizations are inverse
on the images of every $k\in K$.

The image of $s(Q)$ in $G/N$ still maps isomorphically onto $Q$,
because $N\subset K$, and commutes with $K/N$. Thus multiplication
identifies $G/N$ with $C\times Q$. Finally, quotienting by a subgroup
of $[G,G]$ does not change the abelianization: its commutator subgroup
is $[G,G]/N$. Applying this observation gives
\[
 G/[G,G]\simeq (C\times Q)^{\mathrm{ab}}
              \simeq C\oplus Q^{\mathrm{ab}},
\]
as claimed. Each map used here is induced by the original inclusion
or section, so it gives the asserted splitting.
\end{proof}